\documentclass[11pt]{article}

\usepackage[T1]{fontenc}
\usepackage{amsmath,amssymb,amsthm}
\usepackage[a4paper,margin=1in]{geometry}
\usepackage[hidelinks]{hyperref}
\hypersetup{pdftitle={The Pairwise-GCD Sunflower-Free Value at 127},pdfauthor={Xiaojian Zeng}}

\newtheorem{theorem}{Theorem}
\newtheorem{lemma}{Lemma}
\newcommand{\set}[1]{\left\{#1\right\}}

\title{The Pairwise-GCD Sunflower-Free Value at 127}
\author{Xiaojian Zeng\\
PBC School of Finance, Tsinghua University, 100190\\
\texttt{xiaojianzeng@foxmail.com}}
\date{September 6, 2026}

\begin{document}
\maketitle

\begin{abstract}
For a positive integer $N$, let $f_3(N)$ be the maximum size of a subset
$A\subseteq\set{1,\ldots,N}$ containing no three distinct elements whose
three pairwise greatest common divisors are equal.  We prove
$f_3(127)=16$.  The lower bound is an explicit set checked directly over
all $560$ of its triples.  For the upper bound, the existence of a feasible
set of size at least $17$ is encoded by definition-level forbidden-triple
clauses and an exact totalizer cardinality constraint.  Independent
regeneration reproduces the resulting CNF byte for byte, and a DRAT
certificate verifies its unsatisfiability.  The argument is independent of
the previously reported value at $126$.
\end{abstract}

\section{Statement and context}

Write $[N]=\set{1,2,\ldots,N}$.  A set $A\subseteq[N]$ is \emph{feasible}
when there are no pairwise distinct $x,y,z\in A$ such that
\[
 \gcd(x,y)=\gcd(x,z)=\gcd(y,z).
\]
Thus $f_3(N)$ is the maximum cardinality of a feasible subset of $[N]$.
The subscript denotes the forbidden triple and is not an exponent.

The terminology comes from the classical $\Delta$-system, or sunflower,
condition that all pairwise intersections in a family of sets are equal
\cite{ErdosRado1960}.  In the divisor lattice, where the meet is the greatest
common divisor, McKenna defines a sunflower by requiring the gcd of every pair
to equal the gcd of the entire family \cite[Section~2]{McKenna2005}.  For three
integers this is precisely the forbidden configuration above: if their three
pairwise gcds equal $d$, then their common gcd is also $d$.

Zeng has reported $f_3(126)=16$ in a preprint whose stated method also uses
an explicit lower bound and a Boolean formula with a DRAT derivation
\cite{Zeng2026SSRN7395458}.  The accessible bibliographic record and abstract
use the same definition, but the full text and its certificates were not
available for inspection.  Accordingly, we make no priority claim and no
paper-specific assertion beyond that record.  The proof below determines the
adjacent value $N=127$ directly and does not use the result at $126$.

\begin{theorem}\label{thm:main}
The pairwise-GCD sunflower-free value at $127$ is
\[
 f_3(127)=16.
\]
\end{theorem}

\section{Explicit lower bound}

Let
\begin{equation}\label{eq:witness}
 W=\set{14,20,21,24,25,30,36,40,48,63,80,90,96,98,108,125}.
\end{equation}
The entries of $W$ are distinct and lie in $[127]$.  An exact integer
calculation enumerated all $\binom{16}{3}=560$ unordered triples
$a<b<c$ in $W$ and tested whether
\[
 \gcd(a,b)=\gcd(a,c)=\gcd(b,c).
\]
The list of triples passing that test was empty.  The verification script is
only a direct loop over the displayed set and the standard integer gcd
operation; it is included with the computational artifacts.  Hence $W$ is
feasible and
\begin{equation}\label{eq:lower}
 f_3(127)\geq16.
\end{equation}

\section{Certified upper bound}

Introduce Boolean membership variables $x_1,\ldots,x_{127}$.  For every
unordered triple $a<b<c$ satisfying
\begin{equation}\label{eq:bad}
 \gcd(a,b)=\gcd(a,c)=\gcd(b,c),
\end{equation}
include the clause
\begin{equation}\label{eq:clause}
 \neg x_a\lor\neg x_b\lor\neg x_c.
\end{equation}
Exact enumeration of all $\binom{127}{3}=333{,}375$ unordered triples gives
$119{,}176$ clauses of the form \eqref{eq:clause}.  Add a standard totalizer
encoding, using the unary-tree construction of Bailleux and Boufkhad
\cite[Sections~2.2--2.3]{BailleuxBoufkhad2003} through PySAT's
\texttt{CardEnc} interface \cite{IgnatievEtAl2018}, to impose
\begin{equation}\label{eq:card}
 \sum_{i=1}^{127}x_i\geq17.
\end{equation}
The totalizer introduces auxiliary variables and $8{,}890$ clauses.  The
combined formula $F$ has $1{,}015$ variables and $128{,}066$ clauses.

\begin{lemma}\label{lem:encoding}
The formula $F$ is satisfiable if and only if there is a feasible subset of
$[127]$ having at least $17$ elements.
\end{lemma}

\begin{proof}
Given an assignment to the membership variables, the clauses
\eqref{eq:clause} hold exactly when the selected set contains none of the
triples in \eqref{eq:bad}.  Because every unordered triple of distinct
elements of $[127]$ was enumerated, this is exactly feasibility.  The
totalizer has an extension to its auxiliary variables exactly when
\eqref{eq:card} holds.  Consequently an assignment satisfying $F$ selects a
feasible set of size at least $17$, and the indicator vector of any such set
extends to a satisfying assignment of $F$.
\end{proof}

The following exact checks bind this equivalence to the serialized instance.
An audit program independently re-enumerated the forbidden triples, rebuilt
the totalizer, and reproduced the DIMACS file clause for clause and byte for
byte.  Its SHA-256 digest is
\begin{center}
\texttt{8b79ee1c4a921f27dfdb7d4c6d9581ff}\allowbreak
\texttt{98860b85d1caf903b5a500de6763a8c2}.
\end{center}
CaDiCaL 1.9.5 produced a DRAT refutation of that file.  The CaDiCaL solver
family and its proof-output support are described by Biere et al.
\cite{BiereEtAl2024}; the version used here predates their CaDiCaL 2.0 release.
The refutation was checked by \texttt{drat-trim}, the checker for the DRAT
clausal proof format \cite{WetzlerEtAl2014}.  The checker exited successfully
with \texttt{s VERIFIED}; in backward checking it used $114{,}515$ input
clauses and $1{,}794{,}649$ lemmas, comprising $98{,}745{,}797$ resolution
steps.  The DRAT file has SHA-256 digest
\begin{center}
\texttt{57ae38a8f6ff606ef716bdaa488cf2d}\allowbreak
\texttt{93a74c35ec9e16b9437aaac587ba69718}.
\end{center}
Thus the checked certificate proves that $F$ is unsatisfiable.  By
Lemma~\ref{lem:encoding}, no feasible subset of $[127]$ has at least $17$
elements, so
\begin{equation}\label{eq:upper}
 f_3(127)\leq16.
\end{equation}
Equations \eqref{eq:lower} and \eqref{eq:upper} prove
Theorem~\ref{thm:main}.

\section{Reproducibility and limitations}

The supplementary artifacts comprise the witness verifier, the independent
encoding audit, the DIMACS instance, the DRAT trace, and the DRAT checker.
All gcd and enumeration operations are exact.  The upper bound is a certified
finite computation: its validity depends on the definition-to-CNF argument,
the standard totalizer semantics, and the sound checking of the supplied DRAT
trace.  A separate audit also tested the isolated cardinality encoding on
$384$ complete input assignments, covering every selected count from $0$ to
$127$ in three orders; this was an off-by-one check, not a replacement for
the totalizer semantics.

This work determines only $f_3(127)$.  It neither classifies all extremal sets
nor asserts a value at any other $N$.  Because the full text of the closest
identified preprint was unavailable, the literature comparison is necessarily
limited and no novelty or priority conclusion is drawn.

Python programs were used for exact enumeration and independent regeneration;
PySAT 1.8.dev24 supplied the totalizer encoding, CaDiCaL 1.9.5 produced the
unsatisfiability trace, and \texttt{drat-trim} checked it.  A generative AI
system assisted with the computational search, proof audit, literature
screening, and manuscript drafting.  The theorem rests on the explicit
argument and reproducible exact certificates described above; an AI review
pass is not a substitute for external peer review.

\bibliographystyle{plain}
\bibliography{references}

\end{document}
